% Chapter 05: Projections
% Building read models over run traces

\chapter{Projections}
\label{ch:projections}

% ============================================================
\section{What Projections Are}
\label{sec:what-projections}
% ============================================================

A \textbf{projection} is a read-only transformation over a gert run trace. It
consumes the trace event stream (NDJSON format) produced during runbook
execution and produces a domain-specific view: audit logs, timeline summaries,
evidence bundles, compliance reports, etc.

Projections are:

\begin{itemize}
  \item \textbf{Read-only} --- They never modify the trace. They only derive
    views from it.
  \item \textbf{Post-hoc} --- They run \emph{after} execution completes (or
    during execution as a streaming consumer).
  \item \textbf{Domain-specific} --- Different Kits produce different
    projections. A compliance Kit might project an audit log; a workflow Kit
    might project a timeline view.
  \item \textbf{Pure transformations} --- Like the compiler, projections have no
    side effects (no network, no filesystem writes except stdout).
\end{itemize}

% ============================================================
\section{The Trace Event Stream}
\label{sec:trace-events}
% ============================================================

Gert emits a sequence of trace events during execution. Each event is a JSON
object representing a state transition:

\begin{description}
  \item[\texttt{run.started}] --- Execution began.
  \item[\texttt{step.started}] --- A step began executing.
  \item[\texttt{step.completed}] --- A step finished successfully.
  \item[\texttt{step.failed}] --- A step failed.
  \item[\texttt{evidence.captured}] --- Evidence was collected.
  \item[\texttt{approval.requested}] --- Approval gate was reached.
  \item[\texttt{approval.granted}] --- Approval was granted.
  \item[\texttt{approval.rejected}] --- Approval was denied.
  \item[\texttt{variable.set}] --- A variable was bound.
  \item[\texttt{run.completed}] --- Execution finished.
\end{description}

\paragraph{Event format:}

Each event is a JSON object with:

\begin{description}
  \item[\texttt{event}] --- Event type (e.g., \texttt{"step.started"}).
  \item[\texttt{timestamp}] --- ISO 8601 timestamp (UTC).
  \item[\texttt{run\_id}] --- Unique run identifier.
  \item[\texttt{step\_id}] --- Step ID (for step-level events).
  \item[\texttt{payload}] --- Event-specific data (varies by event type).
\end{description}

\paragraph{Example events:}

\begin{minted}{json}
{"event":"run.started","timestamp":"2025-01-15T10:30:00Z","run_id":"r-abc123","payload":{}}
{"event":"step.started","timestamp":"2025-01-15T10:30:01Z","run_id":"r-abc123","step_id":"approve-deploy","payload":{}}
{"event":"evidence.captured","timestamp":"2025-01-15T10:32:00Z","run_id":"r-abc123","step_id":"approve-deploy","payload":{"type":"text","value":"Reviewed deployment checklist. All items OK.","sha256":"e3b0c..."}}
{"event":"approval.granted","timestamp":"2025-01-15T10:32:05Z","run_id":"r-abc123","step_id":"approve-deploy","payload":{"approver":"alice@acme.example"}}
{"event":"step.completed","timestamp":"2025-01-15T10:32:06Z","run_id":"r-abc123","step_id":"approve-deploy","payload":{"exit_code":0}}
{"event":"run.completed","timestamp":"2025-01-15T10:35:00Z","run_id":"r-abc123","payload":{"status":"success"}}
\end{minted}

% ============================================================
\section{Writing a Projection}
\label{sec:writing-projection}
% ============================================================

A projection is a standalone executable that reads NDJSON events on stdin and
writes the projection output to stdout.

\paragraph{Input:} NDJSON stream of trace events (one JSON object per line).

\paragraph{Output:} Projection-specific format (JSON, YAML, Markdown, CSV, etc.).

\paragraph{Exit code:}
\begin{description}
  \item[\texttt{0}] --- Projection succeeded.
  \item[\texttt{1+}] --- Projection failed (malformed input, internal error).
\end{description}

% ============================================================
\section{Use Cases}
\label{sec:projection-use-cases}
% ============================================================

\paragraph{Audit logs:}
Extract all approval events, evidence captures, and governance checks into a
chronological audit log.

\paragraph{Summary reports:}
Produce a high-level summary: which steps ran, how long they took, which
approvers were involved, what evidence was captured.

\paragraph{Compliance evidence bundles:}
Extract all evidence items (text, attachments) and group them by governance
primitive (approval, attestation, redaction). Output as a structured JSON
document for compliance archival.

\paragraph{Timeline views:}
Render a visual timeline of execution (Markdown, HTML, or JSON for frontend
rendering).

\paragraph{Metrics:}
Extract duration metrics, step counts, approval latencies for dashboards.

% ============================================================
\section{The Projection Binary Interface}
\label{sec:projection-interface}
% ============================================================

\paragraph{Invocation:}

\begin{verbatim}
cat trace.ndjson | ./bin/project-audit-log > audit.json
\end{verbatim}

\paragraph{Stdin:} NDJSON trace events (one event per line).

\paragraph{Stdout:} Projection output (format specified by the projection).

\paragraph{Stderr:} Diagnostic messages (optional, not parsed).

% ============================================================
\section{Example: Audit Log Projection}
\label{sec:example-audit-log}
% ============================================================

Let's build a projection that extracts all approval and evidence events into an
audit log.

\paragraph{Input (trace excerpt):}

\begin{verbatim}
{"event":"step.started","timestamp":"2025-01-15T10:30:01Z","run_id":"r-abc123","step_id":"approve-deploy","payload":{}}
{"event":"approval.requested","timestamp":"2025-01-15T10:30:02Z","run_id":"r-abc123","step_id":"approve-deploy","payload":{"approvers":["alice@acme.example","bob@acme.example"]}}
{"event":"evidence.captured","timestamp":"2025-01-15T10:32:00Z","run_id":"r-abc123","step_id":"approve-deploy","payload":{"type":"text","value":"Deployment checklist complete."}}
{"event":"approval.granted","timestamp":"2025-01-15T10:32:05Z","run_id":"r-abc123","step_id":"approve-deploy","payload":{"approver":"alice@acme.example"}}
{"event":"step.completed","timestamp":"2025-01-15T10:32:06Z","run_id":"r-abc123","step_id":"approve-deploy","payload":{}}
\end{verbatim}

\paragraph{Output (audit log JSON):}

\begin{minted}{json}
{
  "run_id": "r-abc123",
  "audit_entries": [
    {
      "timestamp": "2025-01-15T10:30:02Z",
      "step_id": "approve-deploy",
      "event_type": "approval.requested",
      "details": {
        "approvers": ["alice@acme.example", "bob@acme.example"]
      }
    },
    {
      "timestamp": "2025-01-15T10:32:00Z",
      "step_id": "approve-deploy",
      "event_type": "evidence.captured",
      "details": {
        "type": "text",
        "value": "Deployment checklist complete."
      }
    },
    {
      "timestamp": "2025-01-15T10:32:05Z",
      "step_id": "approve-deploy",
      "event_type": "approval.granted",
      "details": {
        "approver": "alice@acme.example"
      }
    }
  ]
}
\end{minted}

% ============================================================
\section{Example: Projection Pseudocode}
\label{sec:projection-pseudocode}
% ============================================================

Here is pseudocode for the audit log projection (Go):

\begin{minted}{go}
package main

import (
  "bufio"
  "encoding/json"
  "os"
)

type TraceEvent struct {
  Event     string                 `json:"event"`
  Timestamp string                 `json:"timestamp"`
  RunID     string                 `json:"run_id"`
  StepID    string                 `json:"step_id,omitempty"`
  Payload   map[string]interface{} `json:"payload"`
}

type AuditEntry struct {
  Timestamp string                 `json:"timestamp"`
  StepID    string                 `json:"step_id"`
  EventType string                 `json:"event_type"`
  Details   map[string]interface{} `json:"details"`
}

type AuditLog struct {
  RunID   string       `json:"run_id"`
  Entries []AuditEntry `json:"audit_entries"`
}

func main() {
  scanner := bufio.NewScanner(os.Stdin)
  var log AuditLog
  var runID string

  for scanner.Scan() {
    var evt TraceEvent
    json.Unmarshal(scanner.Bytes(), &evt)

    if runID == "" {
      runID = evt.RunID
      log.RunID = runID
    }

    // Filter: only audit-relevant events
    if isAuditEvent(evt.Event) {
      log.Entries = append(log.Entries, AuditEntry{
        Timestamp: evt.Timestamp,
        StepID:    evt.StepID,
        EventType: evt.Event,
        Details:   evt.Payload,
      })
    }
  }

  json.NewEncoder(os.Stdout).Encode(log)
}

func isAuditEvent(eventType string) bool {
  auditEvents := []string{
    "approval.requested",
    "approval.granted",
    "approval.rejected",
    "evidence.captured",
  }
  for _, e := range auditEvents {
    if eventType == e {
      return true
    }
  }
  return false
}
\end{minted}

% ============================================================
\section{Projection Declaration in Manifest}
\label{sec:projection-manifest}
% ============================================================

Projections are declared in the Kit manifest:

\begin{minted}{yaml}
projections:
  - name: compliance/audit-log
    description: "Audit log with approvals and evidence"
    output-format: json
  - name: compliance/evidence-bundle
    description: "Evidence bundle for compliance archival"
    output-format: json
  - name: compliance/timeline
    description: "Execution timeline in Markdown"
    output-format: markdown
\end{minted}

\paragraph{Invocation by gert:}

Gert can invoke projections automatically after a run completes:

\begin{verbatim}
gert run runbook.yaml --project compliance/audit-log > audit.json
\end{verbatim}

Or manually:

\begin{verbatim}
cat .gert/traces/r-abc123.ndjson | ./bin/project-audit-log > audit.json
\end{verbatim}

% ============================================================
\section{Streaming Projections}
\label{sec:streaming-projections}
% ============================================================

Projections can run in \textbf{streaming mode}: they consume events as they are
emitted during execution, rather than waiting for the run to complete.

Use cases:

\begin{itemize}
  \item Real-time dashboards (stream events to a web UI).
  \item Incremental compliance logging (write audit entries as they occur).
  \item Alerting (emit alerts when specific events occur).
\end{itemize}

\paragraph{Streaming contract:}

The projection binary should:

\begin{enumerate}
  \item Read events from stdin line by line (NDJSON).
  \item Process each event and emit incremental output to stdout (optional).
  \item Flush stdout after each output (for real-time consumption).
  \item Handle incomplete runs gracefully (if stdin closes before
    \texttt{run.completed}, emit partial output).
\end{enumerate}

Gert can invoke streaming projections with:

\begin{verbatim}
gert run runbook.yaml --stream-to "./bin/project-audit-log" &
\end{verbatim}

The projection runs as a child process, consuming events as they are emitted.
